% Part: lambda-calculus
% Chapter: lambda-definability
% Section: partial-recursive-functions

\documentclass[../../../include/open-logic-section]{subfiles}

\begin{document}

\olfileid{lam}{ldf}{par}
\olsection{الدوال العودية الجزئية \usetoken{S}{lambda definable}}

نحصل على الدوال العودية الجزئية من الدوال الأساس بالتركيب والعودية
البدائية والتصغير غير المحدود. والفرق بينها وبين الدوال العودية العامة
أن الدالة التي يجري عليها البحث غير المحدود لا يشترط أن تكون نظامية.
فإذا لم يوجد موضع ينتهي عنده البحث، لم تكن الدالة الناتجة من التصغير
معرّفة عند تلك الوسائط.

وقد يتبادر أن برهان كون الدوال العودية العامة الكلية !!{lambda definable}
ينتقل بحروفه إلى دعوى أن جميع الدوال العودية الجزئية !!{lambda definable}.
فمثلًا، يكون تركيب~$f$ مع~$g$ !!{lambda defined} بالحد~$\lambd[x][F (G x)]$ متى كانت
$f$ و$g$ !!{lambda defined} بالحدين~$F$ و$G$ على الترتيب. لكن الجزئية تكشف
خللًا في هذا الاستدلال. فكون~$g(x)$ غير معرّفة يمثله أن~$G x$ لا
صورة سوية له. وفي أكثر الأحوال لا تكون لـ~$F (G x)$ صورة سوية كذلك،
وهو الموافق للمطلوب؛ إلا أن ذلك ليس لازمًا. خذ~$F$ مساويًا للحد
$\lambd[x][\lambd[y][y]]$، فإن~$F (G x)$ يختزل إلى الصورة السوية
$\lambd[y][y]$ ولو لم تكن للحد الداخل صورة سوية.

وليس هذا مانعًا نهائيًا. فثمة طرائق بها !!{lambda define} جميع الدوال
العودية الجزئية، مع تمثيل القيم غير المعرفة بحدود لا صورة سوية لها.
إلا أن إنشاءها أدق وأقل ظهورًا من الإنشاء الذي استعملناه للدوال
العودية العامة، فنقتصر هنا على ذكر المبرهنة من غير برهان:

% تصحيح أسلوبي (0381-N1): لا يقال إن النتيجة أُثبتت حيث لا يرد برهان.

\begin{thm}
  جميع الدوال العودية الجزئية !!{lambda definable}.
\end{thm}

\end{document}
